\documentstyle[%


righttag%


,11pt%


,amssymb%


,russian%


,izvru%


]{amsart}





\newcommand{\tl}[3]{\medskip\noindent{\bf#1}\ #2 \dotfill #3 \par} 


\pagestyle{empty}


% \pagestyle{plain}


\begin{document}


% \setcounter{page}{80}


\par


\par


\par





Данный номер является продолжением тематического номера 12 за 2001 г.,


посвященного задачам


оптимального управления и математического программирования.





\bigskip








\centerline{\Large СОДЕРЖАНИЕ}


\bigskip


\bigskip


%%%%%%%%%% Use \newline %%%%%%%%%%%








\tl{Аргучинцева М.А.}


{Вариационная задача об оптимальной форме тела


с минимальным\newline радиационным нагревом поверхности}{3}


\tl{Ащепков Л.Т., Давыдов Д.В.}


{Стабилизация наблюдаемой линейной системы\newline управления


с постоянными интервальными коэффициентами}{11}


\tl{Бадам У.}


{Необходимые условия оптимальности в задачах живучести}{18}


\tl{Габдулхаев Б.Г., Рахимов И.К.}


{Об одном оптимальном сплайн-методе решения\newline операторных


уравнений}{23}


\tl{Исламов Г.Г.}


{О допустимых помехах линейных управляемых систем}{37}


\tl{Коняев Ю.A.}


{Метод унитарных преобразований в теории устойчивости}{41}


\tl{Сазанова Л.А.}


{Устойчивость оптимального синтеза


в задаче быстродействия}{46}


\tl{Ченцов А.Г.}


{К вопросу о построении корректных расширений


в классе конечно-\newline аддитивных мер}{58}


\bigskip


\bigskip





\centerline{КРАТКИЕ~СООБЩЕНИЯ}


\bigskip





\tl{Кокурин М.Ю., Юсупова Н.А.} 


{О необходимых и достаточных условиях медленной\newline


сходимости методов решения линейных некорректных задач}{81}








\vfill


\noindent\raisebox{2pt}{\copyright} Казанский госудаpственный унивеpситет


\end{document}


